% =====================================================================
% 決定力場力学 論文3 — 章立て素案・日本語版(v0.1, 2026-08-16)
% テーマ: 「現実較正」(2026-08-16 原仮定者裁定 — 第121〜122便)。
%   従来この枠にあった「スケール横断の創発」ドラフトは内容不変のまま
%   論文4(dfm-paper4.tex / dfm-paper4-ja.tex)へ繰り下げ。
% 英語版: dfm-paper3.tex(正本 — 章構成・数値は同一)。本稿は骨子のみ。
% 柱(便の時系列順):
%   P1 43″/世紀の力学直接再現(第113〜115便: 実c規約・stateCarry:"double"・比1.002)
%   P2 kF1 安定性の単位系依存(第119便: ループ利得の w=m/d が次元を持つ)
%   P3 kF1 引きずり歳差の較正 — 月の近点回転 8.85年を二体で再現(第120便: D₀ が較正ノブ)
%   P4 共通補正 (D₀=0.006, q=5) — 1組の定数が地球月・水星の両系を同時に満たす
%      (第122便: スピン項の幾何減衰・D₀ 7桁差の解消。tests/exp-kf1c.mjs)
% 「正直な較正」の哲学(全章を貫く): 較正一致 ≠ 機構同定。規約を生んだ失敗
%   (kFrame=0 規約・ソフトニング逆行歳差・窓依存)も結果の一部として報告する。
% =====================================================================
\documentclass[titlepage]{ltjsarticle}

\usepackage{graphicx}
\usepackage{hyperref}
\usepackage{booktabs}

\begin{document}

\title{マッハ的トイ宇宙の現実較正:\\
43″/世紀から月の近点回転まで — 機構主張なしの正直な較正}
\author{(著者ブロックは論文1〜2に同じ)}
\date{2026-08-16(骨子 v0.1)}
\maketitle

% 章構成は英語版 dfm-paper3.tex v0.1 と同一:
% 1. 序論(現実較正 vs 機構同定の定義・失敗を方法とする時系列)
% 2. スケール換算規約(サンプル別指数・実G/実c規約・物理対応ロック Kt=c₀²/G)
% 3. P1: 43″/世紀の力学直接再現(比1.002・ソフトニング逆行歳差の正直な記録)
% 4. P2: 安定性の単位系依存(w=m/d の次元性)
% 5. P3: 運動引きずりによる月の近点回転 8.85年(D₀=0.006+初速−0.32% 共同フィット・
%    実際の主因は太陽摂動=別機構での較正一致)
% 6. P4: 共通補正 (D₀=0.006, q=5)(q*=4.59 の算出・引きずり −5.3e-8=1PNの1/10・
%    Δϖ比 0.988→0.999 純化・被引きずり側スピン無関与の正直な否定結果)
% 7. 環系スコープ(土星 e6+主要衛星6 — タイタン 15.949日 vs 15.945日)
% 8. 較正が許すこと・許さないこと(哲学章)
% 9. Negative claims(論文1/2/4 と同型の固定リスト)
% 10. 再現性(ハーネス・結果JSON・QAゲート・提出時タグ/DOI)

\section*{骨子}
本稿は章立て素案である。本文は英語版確定後に全訳する(論文1〜2の ja 版と同じ運用 —
正本は英語版・相違時は英語版優先)。

\end{document}
